\documentclass[10pt,a4paper]{amsart}
\usepackage{amsmath,amssymb,amsthm,graphicx}
% Retained preamble for the Stacks Project chapter divisors.tex (Divisors).
% The chapter file pinned-original-source.tex carries only \input{preamble} on its first line, so no
% preamble text lives inside the chapter itself.  The Stacks Project uses one single shared preamble for
% every chapter; the snapshot stages it verbatim in sources/stacks-chow/preamble.tex.  The block below is
% copied mechanically, line for line, from that file's lines 14-18 (the Xy-pic setup) and lines 154-217
% (the native theorem declarations together with the mathematical macros).  Nothing was added, removed or
% redefined.  The wrapper class and the amsmath/amssymb/amsthm/graphicx base packages are supplied by the
% assembler, so the class/\externaldocument part of the shared preamble is not repeated here.
% For commutative diagrams we use Xy-pic
\usepackage[all]{xy}

% We use 2cell for 2-commutative diagrams.
\xyoption{2cell}
% Theorem environments.
%
\theoremstyle{plain}
\newtheorem{theorem}[subsection]{Theorem}
\newtheorem{proposition}[subsection]{Proposition}
\newtheorem{lemma}[subsection]{Lemma}

\theoremstyle{definition}
\newtheorem{definition}[subsection]{Definition}
\newtheorem{example}[subsection]{Example}
\newtheorem{exercise}[subsection]{Exercise}
\newtheorem{situation}[subsection]{Situation}

\theoremstyle{remark}
\newtheorem{remark}[subsection]{Remark}
\newtheorem{remarks}[subsection]{Remarks}

\numberwithin{equation}{subsection}

% Macros
%
\def\lim{\mathop{\mathrm{lim}}\nolimits}
\def\colim{\mathop{\mathrm{colim}}\nolimits}
\def\Spec{\mathop{\mathrm{Spec}}}
\def\Hom{\mathop{\mathrm{Hom}}\nolimits}
\def\Ext{\mathop{\mathrm{Ext}}\nolimits}
\def\SheafHom{\mathop{\mathcal{H}\!\mathit{om}}\nolimits}
\def\SheafExt{\mathop{\mathcal{E}\!\mathit{xt}}\nolimits}
\def\Sch{\mathit{Sch}}
\def\Mor{\mathop{\mathrm{Mor}}\nolimits}
\def\Ob{\mathop{\mathrm{Ob}}\nolimits}
\def\Sh{\mathop{\mathit{Sh}}\nolimits}
\def\NL{\mathop{N\!L}\nolimits}
\def\CH{\mathop{\mathrm{CH}}\nolimits}
\def\proetale{{pro\text{-}\acute{e}tale}}
\def\etale{{\acute{e}tale}}
\def\QCoh{\mathit{QCoh}}
\def\Ker{\mathop{\mathrm{Ker}}}
\def\Im{\mathop{\mathrm{Im}}}
\def\Coker{\mathop{\mathrm{Coker}}}
\def\Coim{\mathop{\mathrm{Coim}}}

% Boxtimes
%
\DeclareMathSymbol{\boxtimes}{\mathbin}{AMSa}{"02}

%
% Macros for moduli stacks/spaces
%
\def\QCohstack{\mathcal{QC}\!\mathit{oh}}
\def\Cohstack{\mathcal{C}\!\mathit{oh}}
\def\Spacesstack{\mathcal{S}\!\mathit{paces}}
\def\Quotfunctor{\mathrm{Quot}}
\def\Hilbfunctor{\mathrm{Hilb}}
\def\Curvesstack{\mathcal{C}\!\mathit{urves}}
\def\Polarizedstack{\mathcal{P}\!\mathit{olarized}}
\def\Complexesstack{\mathcal{C}\!\mathit{omplexes}}
% \Pic is the operator that assigns to X its picard group, usage \Pic(X)
% \Picardstack_{X/B} denotes the Picard stack of X over B
% \Picardfunctor_{X/B} denotes the Picard functor of X over B
\def\Pic{\mathop{\mathrm{Pic}}\nolimits}
\def\Picardstack{\mathcal{P}\!\mathit{ic}}
\def\Picardfunctor{\mathrm{Pic}}
\def\Deformationcategory{\mathcal{D}\!\mathit{ef}}

% --- appended 2026-10-12 (Chief): the project own definitions of the reference/slogan/history
% display environments, copied mechanically from the shared Stacks preamble lines 8-12
% (\usepackage{verbatim} plus the three \newenvironment lines, which the chapter preamble keeps
% verbatim from the project preamble). Nothing added, removed or redefined: these environments
% carry no mathematical content and the project itself typesets them as comments. Without these
% three lines a counted statement whose own text uses them fails with Environment slogan undefined.
\usepackage{verbatim}
\newenvironment{reference}{\comment}{\endcomment}
\newenvironment{slogan}{\comment}{\endcomment}
\newenvironment{history}{\comment}{\endcomment}
\usepackage{cleveref}
\crefname{originallemma}{original lemma}{original lemmas}
\Crefname{originallemma}{Original lemma}{Original lemmas}
\makeatletter
\newlabel{lemma-finite-trivialize-invertible-upstairs}{{0BUT}{source}{}{sourcepointer-1}{}}
\newlabel{lemma-finite-trivialize-invertible-upstairs@cref}{{[originallemma][0][]0BUT}{[1][0][]source}{}{}{}}
\makeatother
\crefname{originallemma}{original lemma}{original lemmas}
\Crefname{originallemma}{Original lemma}{Original lemmas}
\makeatletter
\newlabel{lemma-norm-invertible}{{0BCY}{source}{}{sourcepointer-2}{}}
\newlabel{lemma-norm-invertible@cref}{{[originallemma][0][]0BCY}{[1][0][]source}{}{}{}}
\makeatother
\title{A map of invertible sheaves induces a norm map on their norms, and vanishing of the map is detected at a point exactly when the norm map vanishes there}
\author{The Stacks Project Authors}
\date{stacks-project:divisors.tex:local-snapshot-20261012; GNU FDL 1.2 or later}
\begin{document}
\maketitle
\noindent\textit{Editorial note.} All delivered mathematical text is the author's own, copied from the pinned Stacks Project chapter file pinned-original-source.tex; no statement or proof step was rewritten, reordered or removed.\par\medskip
\noindent\textit{Editorial note.} Following the multi-claim convention for this source, source\_url is the source-level URL https://stacks.math.columbia.edu/ and the source label lemma-norm-map-invertible is carried in claim\_key and here rather than in a per-claim tag URL.\par\medskip
\noindent\textit{Editorial note.} The preamble is the chapter's shared Stacks Project preamble, retained as sources/stacks-divisors/retained-preamble.tex: a line-for-line excerpt of the shared Stacks preamble staged with this snapshot. No macro was invented or redefined, and this card adds no preamble text of its own.\par\medskip
\noindent\textit{Editorial note.} Statement span 3695-3715 and proof span 3717-3735 come from the halfA pool card operations/pool-cards/stacks-divisors-halfA.json and were re-read in the pinned chapter file before this card was written.\par\medskip
\noindent\textit{Editorial note.} Reread under the 2026-10-14 prefilter rule: this candidate is 856 bytes over 19 proof lines and the prefilter matched the string 'Some details omitted.' in the last line of the proof. On rereading, that sentence refers to the cocycle verification of the definition just given, while the proof itself supplies the trivialization, the local equations, the definition of the norm map and the reduction of the final equivalence to condition (2) of the definition of a norm map; the candidate was therefore accepted, in the same way as the previously counted card S13-morph-relative-normalization-dvr, and the remaining omission is recorded here rather than counted as a gap in the statement.\par\medskip
\noindent\textit{Editorial note.} Tag self-check per LANE\_BRIEF 10.5 (5): the label lemma-norm-map-invertible resolves in the project's own tag list to canonical tag 0BCZ (line 14297 of the index, full label divisors-lemma-norm-map-invertible). Since reference/public-items.jsonl holds no row whose primary source\_file is divisors.tex, the same-chapter public tag id set is empty and the intersection with the resolved tag id is empty, so the candidate passes.\par\medskip
\noindent\textit{Editorial note.} The gate was run as: python operations/agent-authoring/tools/precard\_gate.py --source divisors.tex --label lemma-norm-map-invertible --start 3695 --end 3715 --proof-start 3717 --proof-end 3735 --local-id 34088, giving RESULT=OK with the label found inside the statement span, disjoint statement and proof spans, exactly one proof environment in the gap between them, distinct statement and proof text, no public interval row for this chapter and occupancy FREE for divisors.tex 3695 3735. Spans were recomputed live by locating the label in the pinned chapter file, and the located text matched the pool byte counts (statement 654 characters, proof 749 characters after whitespace removal).\par\medskip
\noindent\textit{Editorial note.} 2 cross-references occur inside the counted spans; they are listed in external\_references as navigation pointers only.\par\medskip
\noindent\textit{Allowed standard prerequisites (editorial).}\begin{itemize}
\item Invertible sheaves become trivial after a finite pullback along the source of a finite morphism, cited from the same chapter. Cited at proof lines 3718-3721. Produces the open covering of the target over whose preimages the two invertible sheaves are simultaneously trivial, which is what makes the local definition of the norm map possible.
\item Construction of the norm of an invertible sheaf out of a norm of degree d, cited from the same chapter. Cited at statement line 3709 and proof lines 3727-3730. Supplies the target objects of the induced map and the cocycle description whose compatibility has to be verified.
\item Condition (2) in the definition of a norm map of degree d, cited from the same chapter. Cited at proof lines 3732-3733. Converts the pointwise statement about the local equation f\_j into the stated equivalence between vanishing of the map and vanishing of its norm.
\end{itemize}\medskip
\ifdefined\hypertarget\hypertarget{sourcepointer-1}{}\fi
\noindent\textit{Source reference.} Original lemma 0BUT: focused-seeds/AGP/sources/stacks-snapshot-20261011/tags.txt, source lines 14886--14886. Same-chapter reference used in the retained author text at source line(s) 3673, 3721. It is supplied only as a navigation pointer so that the retained text compiles; the counted item delivers the author's own statement and the author's complete proof, and no part of the referenced item is copied in or substituted for the counted argument.\par\medskip
\ifdefined\hypertarget\hypertarget{sourcepointer-2}{}\fi
\noindent\textit{Source reference.} Original lemma 0BCY: focused-seeds/AGP/sources/stacks-snapshot-20261011/tags.txt, source lines 14296--14296. Same-chapter reference used in the retained author text at source line(s) 3709, 3730. It is supplied only as a navigation pointer so that the retained text compiles; the counted item delivers the author's own statement and the author's complete proof, and no part of the referenced item is copied in or substituted for the counted argument.\par\medskip
\begin{lemma}
\label{lemma-norm-map-invertible}
Let $\pi : X \to Y$ be a finite morphism of schemes. Assume there exists
a norm of degree $d$ for $\pi$. For any $\mathcal{O}_X$-linear map
$\varphi : \mathcal{L} \to \mathcal{L}'$
of invertible $\mathcal{O}_X$-modules there is an $\mathcal{O}_Y$-linear
map
$$
\text{Norm}_\pi(\varphi) :
\text{Norm}_\pi(\mathcal{L})
\longrightarrow
\text{Norm}_\pi(\mathcal{L}')
$$
with $\text{Norm}_\pi(\mathcal{L})$, $\text{Norm}_\pi(\mathcal{L}')$
as in Lemma \ref{lemma-norm-invertible}. Moreover, for
$y \in Y$ the following are equivalent
\begin{enumerate}
\item $\varphi$ is zero at a point of $x \in X$ with $\pi(x) = y$, and
\item $\text{Norm}_\pi(\varphi)$ is zero at $y$.
\end{enumerate}
\end{lemma}

\begin{proof}
We choose an open covering $Y = \bigcup V_j$ such that
$\mathcal{L}$ and $\mathcal{L}'$ are trivial over the opens $\pi^{-1}V_j$.
This is possible by
Lemma \ref{lemma-finite-trivialize-invertible-upstairs}.
Choose generating sections
$s_j$ and $s'_j$ of $\mathcal{L}$ and $\mathcal{L}'$
over the opens $\pi^{-1}V_j$. Then $\varphi(s_j) = f_js'_j$
for some $f_j \in \mathcal{O}_X(\pi^{-1}V_j)$.
Define $\text{Norm}_\pi(\varphi)$ to be multiplication
by $\text{Norm}_\pi(f_j)$ on $V_j$. A simple
calculation involving the cocycles used to construct
$\text{Norm}_\pi(\mathcal{L})$, $\text{Norm}_\pi(\mathcal{L}')$
in the proof of Lemma \ref{lemma-norm-invertible}
shows that this defines
a map as stated in the lemma. The final statement follows
from condition (2) in the definition of a norm map of degree $d$.
Some details omitted.
\end{proof}

\end{document}
